\chapter{Conclusion and Future Work}\label{chap:conclAndFutureWork}

\section{Conclusion}

This thesis addressed a major challenge in healthcare systems in Egypt and other developing countries: medical records are often scattered across different healthcare providers. Patient can visit the same doctor but in different hospitals, and each hospital may have its own system for storing medical records, but do not have the data from other institutions. This fragmentation makes it difficult for healthcare professionals to access a complete view of a patient's medical history, which can lead to missed opportunities for early intervention.
Data still is stored on paper, difficult to access, and not easily shared with healthcare providers. These issues make it harder for healthcare professionals to obtain a complete view of a patient's medical history and limit the use of modern AI technologies in healthcare.

To address these challenges, this thesis designed and implemented \textbf{CFI-Care}, a patient-centered healthcare platform that brings medical information together in one secure system while giving patients control over their own data.

The proposed system tackles the key problems identified in this research:

\begin{itemize}
\item \textbf{Fragmented medical records} are addressed through a centralized FHIR-based server that stores and manages healthcare data using international interoperability standards.

```
\item \textbf{Paper-based healthcare records} are addressed through a mobile application that allows patients to scan and upload medical documents directly from their smartphones.

\item \textbf{Limited visualization of patient history} is addressed through an interactive medical history graph that helps healthcare professionals understand diagnoses, treatments, and patient progress over time.

\item \textbf{Privacy and consent concerns} are addressed through secure authentication, role-based access control, and a consent management system that allows patients to control who can access their healthcare information.

\item \textbf{Limited use of AI in healthcare} is addressed through a Retrieval-Augmented Generation (RAG) system that generates responses for doctors and summaries based on verified patient records, improving reliability and reducing human errors and time consumption.
```

\end{itemize}

The main contribution of this work is the integration of these features into a single platform consisting of web, mobile, and backend components. By combining interoperability standards, secure data sharing, patient-controlled consent, healthcare visualization, document digitization, and AI-powered doctor assistance, CFI-Care demonstrates a practical approach to modern healthcare management.

The system was tested ,and the results showed that the different components work correctly together and support the intended functionality of the platform.

\section{Expected Impact}

CFI-Care provides a foundation for improving healthcare record management in Egypt and other developing countries. By using open standards such as FHIR, OAuth 2.0, and OpenID Connect, the platform can be expanded in the future and integrated with other healthcare systems.

The platform also promotes patient involvement by allowing individuals to manage access to their own healthcare information. This can improve trust, transparency, and collaboration between patients and healthcare providers.

Although CFI-Care is currently a prototype and further development is required before large-scale deployment, the project demonstrates that a secure, interoperable, patient-centered, and AI-enabled healthcare system can be successfully implemented using modern technologies. 
The work provides a strong foundation for future research and development in digital healthcare solutions.


\section{Future Work}

\subsection{System Security}

\paragraph{Architecture Hardening and Lifecycle Maturity}
The current implementation provides a robust baseline for secure health data exchange, Meanwhile, 
future iterations will focus on advancing this prototype toward enterprise-grade maturity:

\begin{itemize}
    \item \textbf{Transition to Hardware-Backed Key Orchestration:} While the current LUKS-based storage encryption provides comprehensive at-rest protection, 
    the next phase will integrate an external Key Management Service (KMS) or Hardware Security Module (HSM), 
    which will decouple cryptographic key material from the host file system, enabling programmatic key rotation, tamper-evident audit logging, and ephemeral secret injection.
    \item \textbf{Automated Infrastructure Hardening (IaC):} Move from manual container management to Infrastructure-as-Code (IaC) via hardened Kubernetes (K8s) manifests, to enable the enforcement of \textit{Pod Security Admissions} and ensure least-privilege runtime enforcement.
    \item \textbf{Identity Provider Federation and mTLS:} Scale the identity architecture by transitioning from standard OIDC/OIDC-proxy patterns to mTLS (Mutual TLS) service-to-service authentication; requiring certificate-based validation for all internal API traffic should eliminate the reliance on network-level isolation, establishing a true Zero-Trust architecture.
    \item \textbf{Production Lifecycle Orchestration:} Migrate the current static TLS configuration to automated lifecycle management (using the Let's Encrypt providers) integrated directly into the ingress gateway ensuring consistent, zero-touch certificate rotation across production environments , eliminating manual maintenance and preventing expiry-related downtime.
\end{itemize}

\subsection{Mobile Application}

\subsubsection{Patient adding an Episode of Care}
Currently our system only allows doctors to create an episode of care for a patient. 
In the future, we plan to allow patients to create their own episodes of care, which will enable them to track their health conditions and share their health data with their doctors.
This will empower patients to take a more active role in their healthcare and improve patient engagement.

\paragraph{Technical Implementation}
After an appointment, patients are prompted to upload any requested follow-up medical records, 
such as laboratory results or imaging scans, to the Episode of Care associated with the appointment. 
The Episode of Care is initially created by the physician and serves as a centralized repository for all documents related to the patient's condition. 
This will allow doctors to have a more comprehensive view of the patient's health and make more informed decisions about their care.

% \subsection{unit and widget testing and mockito}
% image picker 
% \subsection{How to monetize the application}
% https://www.udacity.com/enrollment/ud518

% \bibliography{sample}
% \bibliographystyle{ieeetr}
